\chapter{Perron--Frobenius Theory for Channels and Transfer Maps}\label{ch:qpf}

This chapter establishes elementary fixed-point existence for quantum channels
and the Perron--Frobenius theory: existence, positive definiteness, and uniqueness
of fixed points for injective and irreducible transfer maps, canonical
gauge constructions, and the
spectral-radius identification of Perron eigenvalues.
The ergodicity of irreducible channels and the CP-map irreducibility
criterion modeled on Wolf's spectral characterization are developed in the
companion quantum-channel volume~\cite[``Ces\`aro convergence for
irreducible channels'' and ``Irreducibility iff the restricted CP spectral
properties hold'']{QICLean2026Blueprint}.
The conceptual source is Wolf's treatment of irreducible positive
maps~\cite[Chapter~6, especially Theorems~6.2--6.5 and Corollary~6.3]{Wolf2012Quantum}
and~\cite{Evans1978Spectral}.


\section{Positive definiteness}


\begin{theorem}[Transfer-map PSD fixed point --- irreducible version]
    \label{thm:transfer_psd_fp_pd_irred}
    \lean{Kraus.posSemidef_fixedPoint_isPosDef_of_irreducible}
    \leanok
    \uses{def:transfer_map}
    Let $A$ be an MPS tensor whose transfer map $\E_A$ is irreducible. If
    $\rho\geq0$, $\rho\neq0$, and $\E_A(\rho)=\rho$, then $\rho$ is positive
    definite.
\end{theorem}

\begin{proof}
    \leanok
    \uses{thm:transfer_cp}
    The transfer map is completely positive, so
    Theorem~\cite[``PSD fixed point is positive definite --- irreducible version'']{QICLean2026Blueprint} applies.
\end{proof}

\begin{theorem}[PSD fixed point is positive definite]
    \label{thm:psd_fp_pd}
    \lean{Kraus.posSemidef_fixedPoint_isPosDef}
    \leanok
    \uses{def:injective, def:transfer_map}
    Let $A$ be an injective MPS tensor and let $\rho \ge 0$,
    $\rho \neq 0$, be a fixed point of $\E_A$. Then $\rho$ is positive
    definite.
\end{theorem}

\begin{proof}\leanok
    \uses{thm:injective_irreducible, thm:transfer_psd_fp_pd_irred}
    Injectivity implies irreducibility of the transfer map. Apply
    Theorem~\ref{thm:transfer_psd_fp_pd_irred}.
\end{proof}


\section{Uniqueness}


\begin{theorem}[Irreducible fixed-point uniqueness]
    \label{thm:psd_fp_unique_irred}
    \lean{Kraus.posSemidef_fixedPoint_unique_of_irreducible}
    \leanok
    \uses{def:transfer_map}
    Let $A$ be an MPS tensor whose transfer map is irreducible. If
    $\rho,\sigma\geq0$ are fixed points of $\E_A$ and $\rho\neq0$, then
    $\sigma=c\rho$ for some $c\in\C$.
\end{theorem}

\begin{proof}\leanok
    \uses{thm:transfer_cp}
    The transfer map is completely positive, so
    Theorem~\cite[``Irreducible fixed-point uniqueness --- CP version'']{QICLean2026Blueprint} applies.
\end{proof}

\begin{theorem}[Uniqueness of PSD fixed point]
    \label{thm:psd_fp_unique}
    \lean{Kraus.posSemidef_fixedPoint_unique}
    \leanok
    \uses{def:injective, def:transfer_map}
    Let $A$ be an injective MPS tensor. If $\rho,\sigma\geq0$ are nonzero
    fixed points of $\E_A$, then $\sigma=c\rho$ for some $c\in\C$.
\end{theorem}

\begin{proof}\leanok
    \uses{thm:injective_irreducible, thm:psd_fp_unique_irred}
    Injectivity implies irreducibility of the transfer map. Apply
    Theorem~\ref{thm:psd_fp_unique_irred}.
\end{proof}


\section{Existence and the Perron--Frobenius theorem}

\begin{theorem}[Existence of PSD fixed point]\label{thm:psd_fp_exists}
    \lean{Kraus.exists_posSemidef_fixedPoint}
    \leanok
    \uses{def:transfer_map}
    Assume $D \ge 1$. Let $A$ be an MPS tensor with
    $\sum_i (A^i)^\dagger A^i = \Id$, so that $\E_A$ is
    trace-preserving.
    Then there exists $\rho \ge 0$, $\rho \neq 0$,
    with $\E_A(\rho) = \rho$.
\end{theorem}

\begin{proof}
    \leanok
    \uses{thm:transfer_channel}
    Under the normalization $\sum_i (A^i)^\dagger A^i = \Id$,
    the transfer map is a quantum channel
    (Theorem~\ref{thm:transfer_channel}).
    The Ces\`aro fixed-point theorem
    (the corresponding theorem in the companion quantum-channel volume~\cite[``Ces\`aro fixed-point theorem'']{QICLean2026Blueprint}) then provides the
    fixed point; injectivity is not needed for this step.
\end{proof}

\begin{theorem}[Quantum Perron--Frobenius]\label{thm:qpf}
    \lean{Kraus.quantum_perron_frobenius}
    \leanok
    \uses{def:injective, def:transfer_map}
    Assume $D \ge 1$. Let $A$ be an injective MPS tensor with
    $\sum_i (A^i)^\dagger A^i = \Id$.
    Then the transfer map $\E_A$ has a unique positive semidefinite fixed point
    $\rho$ up to scaling, and $\rho$ is positive definite.
\end{theorem}

\begin{proof}
    \leanok
    \uses{thm:psd_fp_exists, thm:psd_fp_pd, thm:psd_fp_unique}
    Combine Theorems~\ref{thm:psd_fp_exists},~\ref{thm:psd_fp_pd},
    and~\ref{thm:psd_fp_unique}.
\end{proof}

\begin{theorem}[Quantum Perron--Frobenius for irreducible transfer maps]
    \label{thm:qpf_irreducible}
    \lean{Kraus.quantum_perron_frobenius_irreducible}
    \leanok
    \uses{def:transfer_map}
    Assume $D\ge 1$. Let $A$ be an MPS tensor such that its transfer map is
    irreducible and $\sum_i (A^i)^\dagger A^i=\Id$, so that $\E_A$ is
    trace-preserving.
    Then $\E_A$ has a unique positive definite fixed point, up to scalar
    multiple.
\end{theorem}

\begin{proof}\leanok
    \uses{thm:psd_fp_exists, thm:transfer_psd_fp_pd_irred, thm:psd_fp_unique_irred}
    Existence is the channel fixed-point theorem
    (Theorem~\ref{thm:psd_fp_exists}). Irreducibility upgrades the fixed point
    to positive definite, and Theorem~\ref{thm:psd_fp_unique_irred} gives
    uniqueness up to scalar multiple.
\end{proof}

\section{Right- and left-canonical gauges}


\begin{definition}[Left-canonical normalization]
    \label{def:left_canonical_tensor}
    \lean{MPSTensor.IsLeftCanonical}
    \leanok
    An MPS tensor $A=(A^i)_i$ is left-canonical when
    \begin{align}
        \sum_i (A^i)^\dagger A^i
         & =\Id.
        \notag
    \end{align}
    Equivalently, its Kraus map is trace-preserving.
\end{definition}

\begin{theorem}[Left-canonical (trace-preserving) gauge]\label{thm:left_canonical_gauge}
    \lean{Kraus.gauged_isTP_of_map_conjTranspose_fixedPoint}
    \lean{gauged_tracePreserving}
    \leanok
    Let $\sigma$ be a fixed point of the adjoint transfer map
    $\sum_i (A^i)^\dagger \sigma A^i = \sigma$,
    and let $S$ be invertible with $S^\dagger S = \sigma$.
    Define $A'^i = SA^iS^{-1}$.
    Then $\sum_i (A'^i)^\dagger A'^i = \Id$, i.e.\ the gauged
    Kraus map is trace-preserving. This is the left-canonical
    normalization.
\end{theorem}

\begin{proof}\leanok
    Conjugating each summand gives
    $(SA^iS^{-1})^\dagger(SA^iS^{-1})
        =(S^\dagger)^{-1}(A^i)^\dagger\sigma A^iS^{-1}$.
    The adjoint fixed-point substitution leaves
    $(S^\dagger)^{-1}\sigma S^{-1}$, and $\sigma=S^\dagger S$ cancels
    the outer factors. See Section~\ref{sec:app_qpf_gauge} for the full
    calculation.
\end{proof}

\begin{theorem}[Square-root trace-preserving gauge from an adjoint fixed point]%
    \label{thm:tp_gauge_from_adjoint_fixed}
    \lean{MPSTensor.tpGauge_isTP_of_transferMap_conjTranspose_fixedPoint}
    \leanok
    \uses{def:transfer_map}
    Let $\sigma$ be positive definite and satisfy
    $\sum_i (A^i)^\dagger \sigma A^i = \sigma$.
    Define
    \begin{align}
        B^i
         & = \sigma^{1/2}A^i\sigma^{-1/2}.
        \label{eq:qpf_tp_gauge}
    \end{align}
    Then $B$ is trace-preserving:
    $\sum_i (B^i)^\dagger B^i = \Id$.
\end{theorem}

\begin{proof}\leanok

    This is the specialization $K_i = A^i$ of
    Theorem~\cite[``Trace-preserving gauge from an adjoint fixed point, Kraus form'']{QICLean2026Blueprint}. Each conjugated
    summand is
    $(B^i)^\dagger B^i
        =\sigma^{-1/2}(A^i)^\dagger\sigma A^i\sigma^{-1/2}$.
    Substituting $\sum_i(A^i)^\dagger\sigma A^i=\sigma$ and cancelling
    the outer inverse square roots gives $\Id$. The complete calculation
    appears in Section~\ref{sec:app_qpf_gauge}.
\end{proof}

\begin{theorem}[Square-root trace-preserving gauge from an adjoint eigenvector]%
    \label{thm:tp_gauge_from_adjoint_eigenvector}
    \lean{MPSTensor.tpGauge_isTP_of_transferMap_conjTranspose_eigenvector}
    \leanok
    \uses{thm:tp_gauge_from_adjoint_fixed}
    Let $\sigma$ be positive definite and satisfy
    $\sum_i (A^i)^\dagger \sigma A^i = r\sigma$ for $r>0$.
    Define
    \begin{align}
        B^i
         & = \sigma^{1/2}(r^{-1/2}A^i)\sigma^{-1/2}.
        \label{eq:qpf_tp_gauge_eigenvector}
    \end{align}
    Then $B$ is trace-preserving:
    $\sum_i (B^i)^\dagger B^i = \Id$.
\end{theorem}

\begin{proof}\leanok
    Rescale $A$ by $r^{-1/2}$. The adjoint transfer map is quadratic in
    the tensor, so the eigenvector equation becomes a fixed-point equation.
    Theorem~\ref{thm:tp_gauge_from_adjoint_fixed} then gives the result.
\end{proof}


\begin{remark}\leanok
    The right-canonical gauge of Theorem~\cite[``Right-canonical (unital) gauge'']{QICLean2026Blueprint}
    and the left-canonical gauge of Theorem~\ref{thm:left_canonical_gauge}
    are generally different similarity transforms: one is built from a fixed point
    of $\E_A$, the other from a fixed point of the adjoint transfer map.
    This section does not claim that a single gauge makes the transfer map
    simultaneously unital and trace-preserving.
\end{remark}


\section{Perron--Frobenius eigenvector existence}%
\label{sec:pf_eigenvector}

This section establishes the existence of a positive definite
eigenvector for the adjoint transfer map of an irreducible MPS
tensor, and uses it to construct a TP-gauge normalization.
The PSD-eigenvector step corresponds to~\cite[Theorem~6.5]{Wolf2012Quantum}
(general positive maps) and~\cite{Evans1978Spectral}; the upgrade to positive
definiteness uses the irreducible-map theory of~\cite[Theorem~6.3]{Wolf2012Quantum}.
The TP-gauge application follows~\cite[Appendix~A]{Cirac2016MPDO_arXiv}.


\begin{theorem}[Irreducibility of the transfer map of an irreducible tensor]%
    \label{thm:transfer_map_irreducible_of_tensor}
    \lean{MPSTensor.isIrreducibleCP_transferMap_of_isIrreducibleTensor}
    \leanok
    \uses{def:irreducible_tensor, def:transfer_map}
    If an MPS tensor $A$ is irreducible, then its transfer map $\E_A$ is
    irreducible.
\end{theorem}

\begin{proof}\leanok
    If $P$ is a nontrivial invariant projection for $\E_A$, then
    \begin{align}
        (\Id-P)A^iP
         & = 0
        \notag
    \end{align}
    for every $i$, which is exactly a nontrivial invariant projection for
    the family $\{A^i\}$, contradicting irreducibility of $A$.
\end{proof}

\begin{theorem}[Positive-definite adjoint eigenvector for irreducible tensors]%
    \label{thm:posDef_adjoint_eigenvector}
    \lean{MPSTensor.exists_posDef_adjoint_eigenvector}
    \leanok
    \uses{def:irreducible_tensor, def:transfer_map}
    Let $A$ be an irreducible MPS tensor with $D > 0$
    and some $A^i \neq 0$.
    Then there exist a positive definite matrix~$\sigma$
    and a positive real $r > 0$ such that
    $\E_A^\dagger(\sigma) = r\sigma$.

    This combines the eigenvector-existence part of
    \cite[Theorem~6.5]{Wolf2012Quantum} with the
    irreducible upgrade to positive definiteness
    (\cite[Theorem~6.3(2)--(3)]{Wolf2012Quantum});
    the application to the adjoint transfer map follows~\cite[Appendix~A]{Cirac2016MPDO_arXiv}.
\end{theorem}

\begin{proof}\leanok
    \uses{thm:transfer_map_irreducible_of_tensor}
    By Theorem~\ref{thm:transfer_map_irreducible_of_tensor}, $\E_A$ is
    irreducible, and since $\mathcal K_A = \E_A$
    (the corresponding theorem in the companion quantum-channel volume~\cite[``Kraus-map form of the transfer map'']{QICLean2026Blueprint}), its Kraus map $\mathcal K_A$ is
    irreducible too (the corresponding theorem in the companion quantum-channel volume~\cite[``Irreducibility in Kraus-map notation'']{QICLean2026Blueprint}). Applying
    Theorem~\cite[``Positive-definite adjoint eigenvector for irreducible Kraus families'']{QICLean2026Blueprint} to $K_i = A^i$ gives
    a positive definite $\sigma$ and a positive real $r > 0$ with
    $\sum_i (A^i)^\dagger \sigma A^i = r\sigma$, i.e.
    $\E_A^\dagger(\sigma) = r\sigma$.
\end{proof}

\begin{theorem}[TP-gauge existence for irreducible tensors]%
    \label{thm:tp_data_irreducible}
    \lean{MPSTensor.exists_tp_data_of_irreducible}
    \leanok
    \uses{def:irreducible_tensor, def:gauge_equiv}
    Let $A$ be an irreducible MPS tensor with $D > 0$
    and some $A^i \neq 0$.
    Then there exist a positive real~$r$, a positive definite matrix~$\sigma$,
    and a tensor~$B$ gauge-equivalent to $r^{-1/2}A$ such that
    \begin{align}
        B^i
         & = \sigma^{1/2}(r^{-1/2}A^i)\sigma^{-1/2},
        \label{eq:qpf_irred_tp_gauge}
    \end{align}
    and $\sum_i (B^i)^\dagger B^i = \Id$.

    This is the ``spectral rescaling $+$ TP gauge'' step of
    \cite[Appendix~A]{Cirac2016MPDO_arXiv}.
    The similarity is generally non-unitary.
\end{theorem}

\begin{proof}\leanok
    \uses{thm:posDef_adjoint_eigenvector, thm:tp_gauge_from_adjoint_eigenvector, lem:tp_gauge_equiv}
    By Theorem~\ref{thm:posDef_adjoint_eigenvector},
    obtain $\sigma$ positive definite and $r > 0$ with
    $\E_A^\dagger(\sigma) = r\sigma$.
    Set $c = r^{-1/2}$ and $A'^i = cA^i$.
    Then $\E_{A'}^\dagger(\sigma) = \sigma$.
    By the square-root trace-preserving eigenvector gauge construction
    (Theorem~\ref{thm:tp_gauge_from_adjoint_eigenvector}), the tensor defined by
    (\ref{eq:qpf_irred_tp_gauge}) satisfies
    $\sum_i (B^i)^\dagger B^i = \Id$.
    Lemma~\ref{lem:tp_gauge_equiv} gives the required gauge equivalence to
    $r^{-1/2}A$.
\end{proof}

\begin{theorem}[Unital-gauge existence for irreducible tensors]%
    \label{thm:unital_data_irreducible}
    \lean{MPSTensor.exists_unital_data_of_irreducible}
    \leanok
    \uses{def:irreducible_tensor}
    Let $A$ be an irreducible MPS tensor with $D > 0$
    and some $A^i \neq 0$.
    Then there exist a positive real~$r$, a positive definite matrix~$\rho$,
    and a tensor~$B$ gauge-equivalent to $r^{-1/2}A$ such that
    \begin{align}
        B^i
         & = r^{-1/2}\rho^{-1/2}A^i\rho^{1/2},
        \label{eq:qpf_irred_unital_gauge}
    \end{align}
    and $\sum_i B^i(B^i)^\dagger = \Id$.

    This is the Perron--Frobenius unital-gauge orientation used in
    \cite[Theorem~4, lines~765--770]{PerezGarcia2007Matrix},
    stated for one irreducible nonzero block.
\end{theorem}

\begin{proof}\leanok
    \uses{thm:posDef_adjoint_eigenvector, thm:unital_gauge_of_pd_fixed_point}
    Apply Theorem~\ref{thm:posDef_adjoint_eigenvector} to the
    conjugate-transposed Kraus family. Irreducibility transfers to that
    family, and the nonzero Kraus hypothesis is preserved by taking adjoints.
    Translating the resulting adjoint eigenvector equation back gives
    $\E_A(\rho)=r\rho$ with $\rho$ positive definite and $r>0$. The unital
    gauge theorem then gives the representative
    (\ref{eq:qpf_irred_unital_gauge}), and the same square-root similarity
    gives gauge equivalence to $r^{-1/2}A$.
\end{proof}
